\chapter{Préface}

Ce livre a une ambition simple~: un seul ouvrage de chimie cohérent,
couvrant tout de la première année d'école à la fin de la licence,
écrit pour des lecteurs du monde entier, d'abord en anglais, et traduit
ici en français.

Le style est concis et rigoureux~: des cours construits à partir de
définitions, d'exemples, de propositions, de méthodes et de schémas, avec
un raisonnement soigneux dès qu'il est accessible au niveau considéré,
suivis d'exercices choisis avec soin. Les résultats énoncés sans
démonstration complète sont explicitement marqués comme admis. Les
solutions complètes de tous les exercices sont rassemblées à la fin du
livre. Les équations chimiques sont toujours ajustées, et chaque nombre
cité pour une substance réelle provient d'une table de référence.

Le cours est réparti en quatre livres. Le premier couvre les douze années
scolaires en un seul volume et n'enseigne jamais deux fois la même
notion~: chaque idée est enseignée une seule fois, l'année où elle trouve
sa place pour la première fois, et un chapitre ultérieur qui en a besoin
s'ouvre sur un court encadré, \emph{Rappel}, avant d'aller plus loin. Les
trois autres livres correspondent aux trois années d'université. Dans le
premier livre, chaque année scolaire est une partie, découpée en
chapitres~; chaque livre universitaire est une année, découpée en
chapitres. Plus l'année est précoce, plus le lecteur visé est jeune~:
davantage d'illustrations, des étapes détaillées plus nombreuses, et des
exercices plus progressifs.

\section*{Comment lire ce livre}

Les énoncés sont numérotés dans chaque chapitre et codés par couleur~:
les définitions en bleu, les théorèmes en rouge, les propositions et
leurs voisins en orange, et les méthodes --- courtes recettes pour les
tâches courantes --- en vert. Les exercices sont notés de ($\star$) pour
les applications directes à ($\star\star\star$) pour les problèmes plus
difficiles. Des encadrés contiennent ce qui ne fait pas partie du
raisonnement~: le rappel de ce qu'un chapitre précédent a enseigné, la
façon dont une manipulation se déroule au laboratoire, une page
d'histoire des sciences, ou les dangers d'un produit.

\section*{Sécurité}

La chimie s'apprend avec de vraies substances, et certaines sont
dangereuses. Les expériences décrites dans ce livre se font dans un
laboratoire scolaire ou universitaire, par un enseignant ou avec lui,
avec les équipements de protection que fournit le laboratoire. Aucune
n'est destinée à être réalisée à la maison.

\section*{Contribuer}

Les sources de ce livre sont dans un dépôt git public~:
\begin{center}
\url{https://github.com/bvirrion/one-chemistry-book}
\end{center}
Les contributions --- nouveaux chapitres, corrections, meilleures
explications, exercices supplémentaires --- sont les bienvenues~; voir
\texttt{CONTRIBUTING.md} à la racine du dépôt.
